\scanpage{49}
\mspage{- 2 -}
Il en résulte que $\operatorname{Ker}(\overset{\handwritten{\sim}}{v})$ est plat \margin{\handwritten{en fait il est \unsure{isom :|isom. à} $\widetilde M$ !}} et que $\overset{\handwritten{\sim}}{M}\handwritten{\times}_S G/\operatorname{Ker}(\overset{\handwritten{\sim}}{v})$ est lisse. Par
ailleurs, vu le choix de $\overset{\handwritten{\sim}}{M}$ , Le groupe $M\times_S G/\overset{\handwritten{\sim}}{M}\handwritten{\times}_S G=M/\overset{\handwritten{\sim}}{M}$ est étale, donc
$M\times_S G/\operatorname{Ker}(\overset{\handwritten{\sim}}{v})$ est lisse , comme extension de groupes lisses. On a évidemment
$\operatorname{Ker}(\overset{\handwritten{\sim}}{v})\handwritten{\subset} P$ , de sorte que $M\times_S G/P$ est \deleted{bien} lisse , comme quotient d'un groupe
lisse.

Démonstration du théorème. L'unicité du prolongement est claire. Comme la
propriété néronienne de $G$ est conservée par passage à l'hensélisé, on peut
supposer $S$ hensélien. Quitte alors à remplacer $\underset{\handwritten{\scriptscriptstyle\sim\!\sim}}{H}$ par sa partie connexe, on
peut , compte tenu du lemme 1 , supposer \deleted{\unsure{$H$}} $\underset{\handwritten{\scriptscriptstyle\sim\!\sim}}{H}$ connexe.

a) Cas où $G$ est semi-stable. Soit $\underset{\handwritten{\scriptscriptstyle\sim\!\sim}}{G}_{\handwritten{\eta}}$ le groupe p-divisible défini par $G_{\handwritten{\eta}}$
et soit $\underset{\handwritten{\scriptscriptstyle\sim\!\sim}}{G}'$ le $S$-groupe p-divisible obtenu à partir du Schéma en groupes $G^{\unsure{o|0|\circ}}$.
% slanted
\margin{\handwritten{``donc'' abusif}}
% «On sait» is circled; the short OK is on its lower left.
\margin{\handwritten{\uline{OK}}}
On sait que $\underset{\handwritten{\scriptscriptstyle\sim\!\sim}}{G}''_{\handwritten{\eta}}=\underset{\handwritten{\scriptscriptstyle\sim\!\sim}}{G}_{\handwritten{\eta}}/\underset{\handwritten{\scriptscriptstyle\sim\!\sim}}{G}'_{\handwritten{\eta}}$ est \inserted{\handwritten{un}} groupe p-divisible \deleted{non ramifié} \inserted{\handwritten{étale}} sur $S$ , \deleted{donc} \handwritten{et}
se prolonge en un $S$-groupe p-divisible , $\underset{\handwritten{\scriptscriptstyle\sim\!\sim}}{G}''$ , étale sur $S$. D'après \uuline{Tate},
le morphisme composé $v_{\handwritten{\eta}}:\underset{\handwritten{\scriptscriptstyle\sim\!\sim}}{H}_{\handwritten{\eta}}\overset{u_{\handwritten{\eta}}}{\handwritten{\longrightarrow}}\underset{\handwritten{\scriptscriptstyle\sim\!\sim}}{G}_{\handwritten{\eta}}\handwritten{\to}\underset{\handwritten{\scriptscriptstyle\sim\!\sim}}{G}''_{\handwritten{\eta}}$ se prolonge en un $S$-morphisme
$v:\underset{\handwritten{\scriptscriptstyle\sim\!\sim}}{H}\handwritten{\to}\underset{\handwritten{\scriptscriptstyle\sim\!\sim}}{G}''$. Comme $\underset{\handwritten{\scriptscriptstyle\sim\!\sim}}{H}$ est connexe et $\underset{\handwritten{\scriptscriptstyle\sim\!\sim}}{G}''$ étale , $v$ est nul. \deleted{P} Il en résulte
que $u_{\handwritten{\eta}}$ se factorise à travers $\underset{\handwritten{\scriptscriptstyle\sim\!\sim}}{G}'_{\handwritten{\eta}}$, donc toujours d'après \uuline{Tate}, se prolonge en un $S$-morphisme $u:\underset{\handwritten{\scriptscriptstyle\sim\!\sim}}{H}\handwritten{\to}\underset{\handwritten{\scriptscriptstyle\sim\!\sim}}{G}'$ c.q.f.d.

% slanted
\margin{\handwritten{? Cas où $G_{\unsure{\eta}}$ est une V.A. seulement !}}
\margin{\handwritten{OK}}
b) Cas général. On sait qu'il existe une \deleted{\unsure{s}}extension de traits $S'\handwritten{\to} S$\handwritten{,}
finie et plate , telle que , si $\handwritten{\eta}'$ est le point générique de $S'$, le $S'$
modèle de Néron $\handwritten{\mathscr G}'$ de $G_{\handwritten{\eta}'}$ soit semi\inserted{\handwritten{-}}stable . D'après a) , le morphisme
$u_{\handwritten{\eta}'}:\underset{\handwritten{\scriptscriptstyle\sim\!\sim}}{H}_{\handwritten{\eta}'}\handwritten{\to}G_{\handwritten{\eta}'}$ se prolonge en un $S'$-morphisme $u':\underset{\handwritten{\scriptscriptstyle\sim\!\sim}}{H}_{S'}\handwritten{\to}\handwritten{\mathscr G}'$ . Rappelons
que le $S$-modèle de Néron de $G_{\handwritten{\eta}}$ , c'est à dire $G$ , s'obtient à partir de $\handwritten{\mathscr G}'$
de la façon suivante (cf. note sur les modèles de Néron prop. 2.5):

On considère le $S$-groupe $S$-néronien $\handwritten{\mathscr G}=\handwritten{\prod}_{S'/S}(\handwritten{\mathscr G}'/S')$. Le groupe $G_{\handwritten{\eta}}$
s'identifie canoniquement à un sous-groupe de $\handwritten{\mathscr G}_{\handwritten{\eta}}$ ; soit $G_1$ l'adhérence
schématique dans $\handwritten{\mathscr G}$ de $G_{\handwritten{\eta}}$ . Alors $G$ est le \textquotedbl lissifié \textquotedbl\ de $G_1$ .
% archive sheet number: 48
